% GENERATED FILE. Edit canonical lessons, then run npm run generate:books.
% canonical-source-hash: fnv1a64:5b5066c4dfd452d8
% canonical-lessons: FA-C386-barabari, FA-C386-nabarabari, FA-C386-edalat, FA-C386-faqr, FA-C386-servat

\chapter{Fairness in society}
\label{ch:fa-fairness-in-society}
\hlchaptermodality{386}

\begin{chapteropening}
\textbf{By the end of this chapter:} \emph{I can talk about equality and inequality, justice, poverty and wealth.}

Complete the last lesson of chapter 386: I can talk about equality and inequality, justice, poverty and wealth.
\end{chapteropening}

\section[\texorpdfstring{\fa{برابری} (barâbari)}{برابری (barâbari)}]{\fa{برابری} (barâbari) — equality}

\label{lesson:FA-C386-barabari}



\begin{quote}
Before the new one: say the Persian for elderly, then the Persian for a religion.
\end{quote}

\subsection*{You'll want to know: \fa{برابری}}
\textbf{\fa{برابری}} (\emph{barâbari}) — \textquotedblleft{}equality\textquotedblright{}.

\begin{quote}
\fa{آنها} \fa{برای} \fa{برابری} \fa{زن} \fa{و} \fa{مرد} \fa{تلاش} \fa{میکنند}.
\end{quote}

\begin{quote}
\emph{ânhâ barâye barâbari-ye zan o mard talâsh mikonand.}
\end{quote}

\begin{quote}
They work for equality between women and men.
\end{quote}

From \emph{barâbar} \textquotedblleft{}equal\textquotedblright{} + \emph{-i}. The opposite is \emph{nâbarâbari}.

\begin{cousinweb}
\textbf{\fa{برابری}} is native Persian, built on \emph{barâbar} \textquotedblleft{}equal, level\textquotedblright{} with the noun suffix \emph{-i}.
\end{cousinweb}

\begin{grammarlens}[title={What you've built}]
One more word for telling, explaining and describing at B1.
\end{grammarlens}

\subsection*{Your turn}
\begin{itemize}
\raggedright
  \item \emph{Say it:} \emph{barâbari}
  \item \emph{Say it:} \emph{barâbari}, once more
  \item \emph{Say it:} a sentence of your own with \emph{barâbari}
  \item \emph{Recall:} say the Persian for elderly, then the Persian for a religion, then say \emph{barâbari} again
\end{itemize}

\begin{tcolorbox}[breakable,colback=teal!4,colframe=teal!35!black,title={Before you move on}]
What does \fa{برابری} mean? (\textquotedblleft{}Equality\textquotedblright{}.) Say it once more.
\end{tcolorbox}
\section[\texorpdfstring{\fa{نابرابری} (nâbarâbari)}{نابرابری (nâbarâbari)}]{\fa{نابرابری} (nâbarâbari) — inequality}

\label{lesson:FA-C386-nabarabari}



\begin{quote}
Before the new one: say the Persian for a religion, then the Persian for equality.
\end{quote}

\subsection*{You'll want to know: \fa{نابرابری}}
\textbf{\fa{نابرابری}} (\emph{nâbarâbari}) — \textquotedblleft{}inequality\textquotedblright{}.

\begin{quote}
\fa{نابرابری} \fa{درآمد} \fa{در} \fa{این} \fa{کشور} \fa{زیاد} \fa{است}.
\end{quote}

\begin{quote}
\emph{nâbarâbari-ye darâmad dar in keshvar ziyâd ast.}
\end{quote}

\begin{quote}
Income inequality in this country is high.
\end{quote}

\emph{Nâ-} \textquotedblleft{}not\textquotedblright{} + \emph{barâbari} \textquotedblleft{}equality\textquotedblright{}.

\begin{cousinweb}
\textbf{\fa{نابرابری}} is native Persian, built from the negative prefix \emph{nâ-} \textquotedblleft{}not\textquotedblright{} + \emph{barâbari} \textquotedblleft{}equality\textquotedblright{}, from \emph{barâbar} \textquotedblleft{}equal\textquotedblright{}.
\end{cousinweb}

\begin{grammarlens}[title={What you've built}]
One more word for telling, explaining and describing at B1.
\end{grammarlens}

\subsection*{Your turn}
\begin{itemize}
\raggedright
  \item \emph{Say it:} \emph{nâbarâbari}
  \item \emph{Say it:} \emph{nâbarâbari}, once more
  \item \emph{Say it:} a sentence of your own with \emph{nâbarâbari}
  \item \emph{Recall:} say the Persian for a religion, then the Persian for equality, then say \emph{nâbarâbari} again
\end{itemize}

\begin{tcolorbox}[breakable,colback=teal!4,colframe=teal!35!black,title={Before you move on}]
What does \fa{نابرابری} mean? (\textquotedblleft{}Inequality\textquotedblright{}.) Say it once more.
\end{tcolorbox}
\section[\texorpdfstring{\fa{عدالت} ('edâlat)}{عدالت ('edâlat)}]{\fa{عدالت} ('edâlat) — justice}

\label{lesson:FA-C386-edalat}



\begin{quote}
Before the new one: say the Persian for equality, then the Persian for inequality.
\end{quote}

\subsection*{You'll want to know: \fa{عدالت}}
\textbf{\fa{عدالت}} (\emph{'edâlat}) — \textquotedblleft{}justice\textquotedblright{}.

\begin{quote}
\fa{مردم} \fa{از} \fa{دولت} \fa{عدالت} \fa{میخواهند}.
\end{quote}

\begin{quote}
\emph{mardom az dowlat 'edâlat mikhâhand.}
\end{quote}

\begin{quote}
People want justice from the government.
\end{quote}

\emph{'Âdel} — fair, righteous. In speech \emph{mikhâhand} is \emph{mikhân}.

\begin{cousinweb}
\textbf{\fa{عدالت}} was borrowed from Arabic \emph{'adāla} \textquotedblleft{}justice, fairness\textquotedblright{} (EtymDB), from the same Arabic family as \emph{'adl} \textquotedblleft{}justice\textquotedblright{}.
\end{cousinweb}

\begin{grammarlens}[title={What you've built}]
One more word for telling, explaining and describing at B1.
\end{grammarlens}

\subsection*{Your turn}
\begin{itemize}
\raggedright
  \item \emph{Say it:} \emph{'edâlat}
  \item \emph{Say it:} \emph{'edâlat}, once more
  \item \emph{Say it:} a sentence of your own with \emph{'edâlat}
  \item \emph{Recall:} say the Persian for equality, then the Persian for inequality, then say \emph{'edâlat} again
\end{itemize}

\begin{tcolorbox}[breakable,colback=teal!4,colframe=teal!35!black,title={Before you move on}]
What does \fa{عدالت} mean? (\textquotedblleft{}Justice\textquotedblright{}.) Say it once more.
\end{tcolorbox}
\section[\texorpdfstring{\fa{فقر} (faqr)}{فقر (faqr)}]{\fa{فقر} (faqr) — poverty}

\label{lesson:FA-C386-faqr}



\begin{quote}
Before the new one: say the Persian for inequality, then the Persian for justice.
\end{quote}

\subsection*{You'll want to know: \fa{فقر}}
\textbf{\fa{فقر}} (\emph{faqr}) — \textquotedblleft{}poverty\textquotedblright{}.

\begin{quote}
\fa{فقر} \fa{هنوز} \fa{یک} \fa{مشکل} \fa{بزرگ} \fa{است}.
\end{quote}

\begin{quote}
\emph{faqr hanuz yek moshkel-e bozorg ast.}
\end{quote}

\begin{quote}
Poverty is still a big problem.
\end{quote}

\emph{Faqir} \textquotedblleft{}poor\textquotedblright{} belongs to the same family.

\begin{cousinweb}
\textbf{\fa{فقر}} was borrowed from Arabic \emph{faqr} \textquotedblleft{}poverty\textquotedblright{} (EtymDB), from the same Arabic family as \emph{faqīr} \textquotedblleft{}poor\textquotedblright{}, the source of Persian \emph{faqir}.
\end{cousinweb}

\begin{grammarlens}[title={What you've built}]
One more word for telling, explaining and describing at B1.
\end{grammarlens}

\subsection*{Your turn}
\begin{itemize}
\raggedright
  \item \emph{Say it:} \emph{faqr}
  \item \emph{Say it:} \emph{faqr}, once more
  \item \emph{Say it:} a sentence of your own with \emph{faqr}
  \item \emph{Recall:} say the Persian for inequality, then the Persian for justice, then say \emph{faqr} again
\end{itemize}

\begin{tcolorbox}[breakable,colback=teal!4,colframe=teal!35!black,title={Before you move on}]
What does \fa{فقر} mean? (\textquotedblleft{}Poverty\textquotedblright{}.) Say it once more.
\end{tcolorbox}
\section[\texorpdfstring{\fa{ثروت} (servat)}{ثروت (servat)}]{\fa{ثروت} (servat) — wealth}

\label{lesson:FA-C386-servat}



\begin{quote}
Before the new one: say the Persian for justice, then the Persian for poverty.
\end{quote}

\subsection*{You'll want to know: \fa{ثروت}}
\textbf{\fa{ثروت}} (\emph{servat}) — \textquotedblleft{}wealth\textquotedblright{}.

\begin{quote}
\fa{ثروت} \fa{او} \fa{از} \fa{پدرش} \fa{به} \fa{او} \fa{رسیده} \fa{است}.
\end{quote}

\begin{quote}
\emph{servat-e u az pedarash be u reside ast.}
\end{quote}

\begin{quote}
His wealth came to him from his father.
\end{quote}

\emph{Servatmand} — wealthy. \fa{ث} sounds like s.

\begin{cousinweb}
\textbf{\fa{ثروت}} was borrowed from Arabic \emph{tharwa} \textquotedblleft{}wealth\textquotedblright{} (EtymDB); Persian says the Arabic \emph{th} as \emph{s}.
\end{cousinweb}

\begin{grammarlens}[title={What you've built}]
One more word for telling, explaining and describing at B1.
\end{grammarlens}

\subsection*{Your turn}
\begin{itemize}
\raggedright
  \item \emph{Say it:} \emph{servat}
  \item \emph{Say it:} \emph{servat}, once more
  \item \emph{Say it:} a sentence of your own with \emph{servat}
  \item \emph{Recall:} say the Persian for justice, then the Persian for poverty, then say \emph{servat} again
\end{itemize}

\begin{tcolorbox}[breakable,colback=teal!4,colframe=teal!35!black,title={Before you move on}]
What does \fa{ثروت} mean? (\textquotedblleft{}Wealth\textquotedblright{}.) Say it once more.
\end{tcolorbox}
